% ============================================================================
% 2.6 主程序：把一切串起来
% ============================================================================
\subsection{hello.s：主程序入口}

主程序的工作非常简单——初始化堆之后，就进入一个循环：
读一行→复制→存到平行数组→打印→找空格。

\begin{lstlisting}[style=asm,caption=hello.s — 主程序]
.text
.global _main
.extern exit
.extern memory_init
.extern memory_alloc
.extern memory_free
.extern sys_write
.extern sys_read
.extern test_write
.extern test_memory
.extern test_read
.extern test_write_memory
.extern store_string_ptr_len
.extern get_string_ptr
.extern get_string_len
.extern store_ptr_done
.extern print_index_strings
.extern find_two_spaces

_main:
    bl memory_init             ; 初始化堆
    mov x19, #0               ; x19 = 循环计数器（也是字符串序号）

loop:
    mov x20, #0
    bl store_string           ; 读一行，分配内存，存入数组

    mov x20, #0
    bl get_string_ptr_len_spaces ; 顺便检查空格位置

    add x19, x19, #1
    cmp x19, #1
    b.eq do_exit
    b loop

do_exit:
    b exit
\end{lstlisting}

\subsection{store\_string：读入并保存一行}

\begin{lstlisting}[style=asm,caption=hello.s — store\_string]
store_string:
    stp x29, x30, [sp, #-16]!
    bl test_read              ; 把一行
    bl copy_string            ; copy_string(x0) 从 input_buf 复制到新分配的堆内存
    mov x1, x20                ; 序号
    bl store_string_ptr_len     ; 存入平行数组
    ldp x29, x30, [sp], #16
    ret
\end{lstlisting}

\subsection{get\_string\_ptr\_len\_spaces：打印并解析空格}

\begin{lstlisting}[style=asm,caption=hello.s — 打印+解析]
get_string_ptr_len_spaces:
    stp x29, x30, [sp, #-16]!

    bl print_string           ; 把刚存入的字符串再取出来打印
    stp x1, x2, [sp, #-16]! ; 保存打印返回的指针和长度
    bl find_two_spaces        ; 然后扫描字符串找空格位置
    ldp x2, x3, [sp], #16

get_string_ptr_len_spaces_done:
    ldp x29, x30, [sp], #16
    ret
\end{lstlisting}

\subsection{理解整体数据流}

整个程序的数据流如下：

\begin{verbatim}
stdin → sys_read → input_buf → copy_string(从 input_buf 到 heap)
        → (ptr, len) → store_string_ptr_len
        → print_index_strings(x19) → sys_write → stdout
        → find_two_spaces(buffer) → x0=第一个空格, x1=第二个空格
\end{verbatim}

程序循环一次（\texttt{x19} 递增一，直到 \texttt{x19 = 1} 时退出。
你可以把 \texttt{cmp x19, \#1} 改成更大的数字来延长程序执行更多轮。

\subsection{主循环限制}

这里的循环在 \texttt{x19} 计数，
\textbf{最多能存 MAX\_STRINGS（32）条字符串}。超过会被
\texttt{store\_string\_ptr\_len} 里的 \texttt{cmp x1, \#MAX\_STRINGS; b.hs store\_ptr\_done} 拦截掉——如果 \texttt{x19} 超过 32，函数什么都不做就返回。
